\documentclass[11pt]{article}
\usepackage[margin=1in]{geometry}
\usepackage{amsmath,amssymb,hyperref,listings}
\title{Supervised ConvNeXt PAR for UPAR Track 2}
\author{}
\date{September 2026}
\begin{document}
\maketitle

\begin{abstract}
The current strongest direction for UPAR Track 2 is supervised pedestrian attribute recognition (PAR), not zero-shot text retrieval or the current hyperbolic prototype model. A ConvNeXt image encoder predicts the 40 UPAR attributes directly. Retrieval compares the predicted attribute probabilities of each gallery image with a binary attribute query using calibrated, attribute-weighted L1 distance.
\end{abstract}

\section{Problem}
Each gallery image $x_i$ has a binary attribute vector $a_i\in\{0,1\}^{40}$ and a semantic query ID $s_i$. A query is also a binary vector:
\[
q\in\{0,1\}^{40}.
\]
The system returns a distance matrix $D\in\mathbb{R}^{Q\times G}$, where lower values rank images earlier. The primary metric is mADM; mAP and Rank-1 are secondary metrics.

\section{Attributes are not text}
UPAR queries are structured labels, not natural-language sentences. For example, a query vector may activate the attributes Young and Gender-Female. The current model therefore does not use a text encoder. It directly consumes the fixed 40-dimensional vector.

This is deliberate: zero-shot CLIP and Hyper3-CLIP were trained on general image-text data, while UPAR attributes have dataset-specific definitions and annotation conventions. The experiments showed that supervised attribute prediction is substantially stronger than zero-shot prompt matching.

\section{ConvNeXt image encoder}
The image encoder is an ImageNet-pretrained ConvNeXt:
\[
f_i=E_\theta(x_i),\qquad f_i\in\mathbb{R}^{d_f}.
\]
ConvNeXt-Tiny produces a 768-dimensional pooled feature. Small and Base produce larger features. The current main experiment starts with Tiny because it gave the best capacity-to-generalization tradeoff in the available results.

The backbone uses convolutional stages, depthwise convolutions, pointwise projections, LayerNorm, and global average pooling. ConvNeXt uses LayerNorm rather than BatchNorm. LayerNorm has no batch running statistics, so small physical batches and gradient accumulation do not create the usual BatchNorm problem. Replacing pretrained LayerNorm with GroupNorm is unnecessary and would change the pretrained architecture.

\section{40-way attribute head}
The classifier is a dropout layer followed by one linear head:
\[
z_i=Wf_i+b,\qquad z_i\in\mathbb{R}^{40}.
\]
Each logit corresponds to one fixed UPAR attribute. The predicted probability is
\[
p_{ik}=\sigma(z_{ik})=\frac{1}{1+\exp(-z_{ik})}.
\]
The gallery representation is therefore
\[
p_i=(p_{i1},\ldots,p_{i40})\in[0,1]^{40}.
\]
No text embedding is required.

\section{Focal multi-label loss}
The base multi-label objective is binary cross-entropy with logits:
\[
\mathcal{L}_{\mathrm{BCE}}=-\frac{1}{40B}\sum_{i=1}^{B}\sum_{k=1}^{40}
\left[a_{ik}\log\sigma(z_{ik})+(1-a_{ik})\log(1-\sigma(z_{ik}))\right].
\]
UPAR attributes are imbalanced, so the trainer uses focal BCE. Define
\[
p_{t,ik}=a_{ik}\sigma(z_{ik})+(1-a_{ik})(1-\sigma(z_{ik})).
\]
Then
\[
\mathrm{FL}_{ik}=\alpha(1-p_{t,ik})^\gamma\mathrm{BCE}_{ik},
\qquad
\mathcal{L}_{\mathrm{focal}}=\frac{1}{40B}\sum_{i,k}\mathrm{FL}_{ik}.
\]
The current settings use $\gamma=2$ and label smoothing $\epsilon=0.05$:
\[
\tilde a=(1-\epsilon)a+\frac{\epsilon}{2}.
\]
Focal loss reduces the influence of easy, common attributes and emphasizes difficult examples.

\section{EMA and checkpointing}
An exponential moving average is updated after optimizer steps:
\[
\theta_{\mathrm{EMA}}=\rho\theta_{\mathrm{EMA}}+(1-\rho)\theta.
\]
The EMA model is used for validation and checkpointing. Training loss alone is not the retrieval metric; each candidate must be evaluated with mADM, mAP, and Rank-1.

\section{Attribute retrieval distance}
For binary query $q$ and predicted gallery probabilities $p_i$, the unweighted distance is
\[
d_{\mathrm{L1}}(q,p_i)=\sum_{k=1}^{40}|q_k-p_{ik}|.
\]
Because $q_k$ is binary,
\[
|q_k-p_{ik}|=p_{ik}+q_k(1-2p_{ik}),
\]
so the full distance matrix can be computed as
\[
d_{\mathrm{L1}}(q,p_i)=\sum_kp_{ik}+\sum_kq_k(1-2p_{ik}).
\]

The evaluator uses attribute weights:
\[
d_w(q,p_i)=\sum_{k=1}^{40}w_k|q_k-p_{ik}|,
\]
with the equivalent memory-safe formula
\[
d_w(q,p_i)=\sum_kw_kp_{ik}+\sum_kq_kw_k(1-2p_{ik}).
\]
This avoids allocating a $Q\times G\times40$ tensor during validation.

\section{Calibration and weighting}
Calibration and retrieval weights must be fitted on the training split and then applied to validation or test. Fitting them on validation labels would leak evaluation information and produce an optimistic score.

The evaluator estimates error weights from attribute prediction errors and distribution-based weights from gallery frequency. These are combined into $w_k$. The clean protocol is:
\[
\texttt{--calib-split train},\qquad\texttt{--split val}.
\]

\section{mADM}
The ground-truth degree of match for query $q_j$ and gallery labels $a_i$ is
\[
\mathrm{DoM}(q_j,a_i)=1-\frac{1}{40}\sum_{k=1}^{40}|q_{jk}-a_{ik}|.
\]
The ranking should place high-DoM images near the top, while exact semantic-ID matches provide the retrieval relevance signal. This is why improving the 40 attribute probabilities directly is more aligned with mADM than generic image-text similarity.

\section{8 GB training}
With physical batch $B$ and accumulation factor $A$, each loss is divided by $A$ before backpropagation:
\[
\widehat{\mathcal{L}}_t=\frac{\mathcal{L}_t}{A}.
\]
Gradients are accumulated for $A$ batches, clipped once, and used for one optimizer update. The approximate effective batch is
\[
B_{\mathrm{effective}}=B\times A.
\]
AMP reduces activation memory. The Tiny one-epoch probe demonstrated that physical batch 64 fits on the 8 GB GPU, so the current full run uses $B=64$ and $A=1$. If it fails, use $B=32$ and $A=2$.

\section{Recommended experiment}
\begin{lstlisting}[basicstyle=\ttfamily\small,breaklines=true]
python -m hyppar.scripts.train_par \\
  --name par_convnext_tiny_30 \\
  --backbone convnext_tiny \\
  --epochs 30 \\
  --batch-size 64 \\
  --grad-accum-steps 1 \\
  --lr 1e-4 \\
  --dropout 0.3 \\
  --focal-gamma 2 \\
  --label-smoothing 0.05 \\
  --amp \\
  --num-workers 2
\end{lstlisting}

Leakage-free evaluation:
\begin{lstlisting}[basicstyle=\ttfamily\small,breaklines=true]
python -m hyppar.scripts.evaluate_par \\
  --ckpt checkpoints/par_convnext_tiny_30/model.pt \\
  --calib-split train \\
  --split val \\
  --batch-size 64 \\
  --num-workers 2 \\
  --save-artifacts \\
  --out checkpoints/par_convnext_tiny_30/eval_val_par_clean.json
\end{lstlisting}

Packaging:
\begin{lstlisting}[basicstyle=\ttfamily\small,breaklines=true]
python -m hyppar.scripts.package_par_submission \\
  --ckpt checkpoints/par_convnext_tiny_30/model.pt \\
  --out checkpoints/par_task2_par_convnext_tiny_30.zip
\end{lstlisting}

\section{Current evidence}
The one-epoch supervised Tiny probe reached approximately $0.524$ mADM with training-split calibration. This is substantially stronger than the current hyperbolic and zero-shot CLIP experiments. The full experiment tests whether longer supervised PAR training, calibration, and weighted attribute retrieval can approach $0.60$ or higher.

\end{document}
